% ---
% id: meat-core/lem:branch-support
% kind: lem
% title: "Properties of the maximal branch support"
% summary: >-
%   Every branch operator the program forms is supported in the maximal branch
%   support, with equality at a positive definite Choi block; the support is
%   the range of the stacked Kraus matrix. It is always smaller than the
%   ambient space and in general not invariant under the key pinching.
% prominence: standard
% status: implemented
% uses: [def:branch-support, def:branch, lem:admissible-alice-subspace, def:untrusted-choi, lem:choi-marginal, def:admissible-declaration]
% mentions: [thm:gradient, meat-core/frank-wolfe/thm:firstorder, meat-core/frank-wolfe/req:support-rank-check]
% assumes: []
% refs: []
% provenance: newly-derived
% implements: certiqs_key.meat_core.objective.gradient:is_support_maximal
% lean:
% ---

Four facts about the \tref{maximal branch support} are used by the gradient
of Theorem~\ref{thm:gradient} and by the first-order argument of the
Frank--Wolfe submodule: it bounds the support of every branch operator the
program forms, it is the range of the stacked Kraus matrix, it is always
strictly smaller than the ambient space, and it is in general not invariant
under the \tref{key pinching}. The bound in~(ii) is not a number a run computes:
$\operatorname{rank}\tau_A$ enters it only as a mathematical quantity.
Hypothesis (H-SUPP) of Theorem~\ref{thm:firstorder} compares a rank with
$\dim W_{\hat c,s}$, and how that comparison is decided is
Requirement~\ref{req:support-rank-check}.

\begin{lemma}\label{lem:branch-support}
Let $\hat c\in\Chkey$ and $s\in\SecSet$, and let $W_{\hat c,s}$ and
$\mathbf K_{\hat c,s}$ be as in Definition~\ref{def:branch-support}.
\begin{enumerate}[label=(\roman*),leftmargin=2.2em]
\item In trusted-source mode every $\omega_s=\chi_{\Qsys_s}(\Delta)$ with
  $\Delta\in\Pos{A_0\Qsys_s}$, in particular every feasible $\omega_s$
  and the image of the positive and the negative part of every Hermitian
  direction, satisfies
  $\operatorname{supp}\omega_s
  =(W_{\tau}\otimes\idon{})(\operatorname{supp}\Delta)
  \subseteq\operatorname{supp}\tau_A\otimes\Hs{\Qsys_s}
  =\Sadm\otimes\Hs{\Qsys_s}$; in untrusted-source mode
  $\Sadm=\Hs{A}$ and the inclusion holds for every
  $\omega_s\in\Pos{A\Qsys_s}$. In either mode, therefore,
  \begin{equation}\label{eq:Wcs-bound}
    \operatorname{supp}\cG_{\hat c,s}(\omega_s)\subseteq W_{\hat c,s},
  \end{equation}
  with equality whenever $\omega_s$ attains that maximal support,
  $\operatorname{supp}\omega_s=\Sadm\otimes\Hs{\Qsys_s}$. In
  trusted-source mode $J_s\succ0$ for the Choi block suffices for it, and
  on a tight support frame is equivalent to it; in untrusted-source mode it
  is $\omega_s\succ0$ on $\Hs{A\Qsys_s}$, which is $J_s\succ0$, as
  $J_s=\omega_s$ there (Definition~\ref{def:untrusted-choi}\,(a)).
\item $W_{\hat c,s}$ is the range of one matrix, and its dimension is
  bounded:
  \begin{equation}\label{eq:Kstack-rank}
    W_{\hat c,s}=\operatorname{ran}\mathbf K_{\hat c,s},
    \qquad
    \dim W_{\hat c,s}=\operatorname{rank}\mathbf K_{\hat c,s}
    \;\le\;\big|\Kann_{\mathrm{pass}}(\hat c)\big|\cdot
    \dim\Sadm\cdot d_s ,
  \end{equation}
  where $\dim\Sadm=\operatorname{rank}\tau_A$ in trusted-source mode and
  $d_A$ in untrusted-source mode.
\item $\dim W_{\hat c,s}$ is \emph{strictly} smaller than
  $\dim\Hs{R\,XY\,C_AC_B\,A\Qsys_s}$, so $0$ always lies in the ambient
  spectrum of a branch operator $X_{\hat c,s}=\cG_{\hat c,s}(\omega_s)$ of~(i)
  and ``maximal support'' is never the statement that $X_{\hat c,s}\succ0$
  on the ambient space.
\item $W_{\hat c,s}$ is \emph{not} in general $\cZ$-invariant.
\end{enumerate}
\end{lemma}

\begin{proof}
\emph{Support of a sum.} For $\varrho\succeq0$ the support of
$\sum_{\kappa}K_{\kappa}\varrho K_{\kappa}^{\dagger}$ is
$\sum_{\kappa}K_{\kappa}(\operatorname{supp}\varrho)$: a vector $\ket w$
in the kernel of the sum has
$\bra wK_{\kappa}\varrho K_{\kappa}^{\dagger}\ket w=0$, hence
$\varrho^{1/2}K^{\dagger}_{\kappa}\ket w=0$, for every $\kappa$, and
conversely each summand has range in $K_{\kappa}(\operatorname{supp}
\varrho)$.

(i) In trusted-source mode,
$\operatorname{supp}\omega_s=(W_{\tau}\otimes\idon{})(\operatorname{supp}\Delta)$
by the support of a sum with the single operator $W_{\tau}\otimes\idon{}$;
it lies in $\operatorname{ran}W_{\tau}\otimes\Hs{\Qsys_s}$, and
$\operatorname{ran}W_{\tau}=\operatorname{supp}\tau_A$
(Lemma~\ref{lem:choi-marginal}\,(iii)), which is $\Sadm$ by \eqref{eq:SA}. In
untrusted-source mode $\Sadm=\Hs{A}$ and the inclusion is trivial. The
support of a sum then gives
$\operatorname{supp}\cG_{\hat c,s}(\omega_s)
=\sum_{\kappa}K^{(s)}_{\kappa}(\operatorname{supp}\omega_s)
\subseteq\sum_{\kappa}K^{(s)}_{\kappa}(\Sadm\otimes\Hs{\Qsys_s})
=W_{\hat c,s}$, with equality where
$\operatorname{supp}\omega_s=\Sadm\otimes\Hs{\Qsys_s}$. That is the case
when $J_s\succ0$, since then
$\operatorname{supp}\omega_s=\operatorname{ran}W_{\tau}\otimes\Hs{\Qsys_s}$,
and when $\omega_s\succ0$ on $\Hs{A\Qsys_s}$ in untrusted-source mode. On a
tight frame the converse holds as well: $W_{\tau}$ is then injective
(Lemma~\ref{lem:choi-marginal}\,(iii)), so $W_{\tau}\otimes\idon{}$ maps
$\operatorname{supp}J_s$ injectively onto $\operatorname{supp}\omega_s$,
and $\operatorname{supp}\omega_s=\operatorname{ran}W_{\tau}\otimes
\Hs{\Qsys_s}$, of dimension $d_{A_0}d_s$, forces
$\operatorname{supp}J_s=\Hs{A_0\Qsys_s}$, that is $J_s\succ0$.

(ii) The range of a side-by-side listing is the sum of the ranges of its
blocks, and
$\operatorname{ran}(W_{\tau}\otimes\idon{\Qsys_s})
=\operatorname{ran}W_{\tau}\otimes\Hs{\Qsys_s}
=\Sadm\otimes\Hs{\Qsys_s}$ by
$\operatorname{ran}W_{\tau}=\operatorname{supp}\tau_A$
(Lemma~\ref{lem:choi-marginal}\,(iii)), so
$\operatorname{ran}\mathbf K_{\hat c,s}
=\sum_{\kappa}K^{(s)}_{\kappa}(\Sadm\otimes\Hs{\Qsys_s})$, which is
\eqref{eq:Wcs}. The rank of a side-by-side listing is at most the sum of
the ranks of its blocks, and each block has rank at most
$\operatorname{rank}(W_{\tau}\otimes\idon{\Qsys_s})=\dim\Sadm\cdot d_s$,
which gives the inequality. In untrusted-source mode the same two steps
apply with $\idon{A\Qsys_s}$ in place of $W_{\tau}\otimes\idon{\Qsys_s}$.

(iii) For a \tref{pass class} $\kappa$ of $\hat c$ every
$(x,v)\in a^{-1}(\kappa)$ has $r(x,v)\in\Zval$, by announcement-determined
sifting and Definition~\ref{def:admissible-declaration}\,(iv), (v); so every Kraus operator
$K^{(s)}_{\kappa}$ of \eqref{eq:Wcs} carries an $R$-ket $\ket{j}_R$ with
$j\in\Zval$, and $W_{\hat c,s}\subseteq\operatorname{span}\{\ket{j}_R:
j\in\Zval\}\otimes\Hs{XY\,C_AC_B\,A\Qsys_s}$. The nonzero vectors
$\ket{\perp}_R\otimes\xi$ are orthogonal to that subspace, so
$\dim W_{\hat c,s}$ is smaller than the ambient dimension. By~(i),
$\operatorname{supp}X_{\hat c,s}\subseteq W_{\hat c,s}$, so those vectors
lie in the ambient kernel of $X_{\hat c,s}$, and $0$ is an eigenvalue of
$X_{\hat c,s}$ on the ambient space.

(iv) The declaration is free to let the raw-key value $r(x,v)$ vary with
$x$ inside one announcement class, since
Definition~\ref{def:admissible-declaration}\,(viii) constrains only that the key holder can
compute it, and then a single Kraus operator \eqref{eq:kraus} has range
vectors spread over several $R$-eigenvalues. $Z_j$ applied to such a
vector $K^{(s)}_{\kappa}\ket\psi$ retains only the terms with
$r(x,v)=j$, which is in general not $K^{(s)}_{\kappa}\ket{\psi'}$ for any
$\ket{\psi'}$, nor in the range of any other Kraus operator (the classes
$\kappa$ occupy orthogonal $C_AC_B$ sectors). An instance: $d_A=1$, one
Alice outcome $x$, $\Qsys=\mathbb{C}$ with one sector $s$, two outcomes
$v_0,v_1$ with $F_{v_0}=F_{v_1}=\frac12$, both outcome pairs in one
announcement class $\kappa$, which passes sifting, and the key held by
Bob with $r(x,v_0)=0$ and $r(x,v_1)=1$, a function of his readout; let
$\hat c\coloneqq\phi(\kappa)$ have no other pass class. Then
$W_{\hat c,s}$ is the line spanned by
$\big(\ket{0}_R\ket{x,v_0}+\ket{1}_R\ket{x,v_1}\big)\ket{\kappa}$, and
$Z_0$ maps that vector to $\ket{0}_R\ket{x,v_0}\ket{\kappa}$, which is not
on the line.
\end{proof}
